\begin{recipe}
[% 
    preparationtime = {\unit[1]{h}\unit[30]{min} + \unit[24]{h} Kühlzeit + \unit[24]{h} Marinierzeit},
    bakingtime = {\unit[75]{min} Dämpfen + \unit[15]{min} Anbraten},
    portion = {\portion{6}},
    source = {\protect\recipesource{https://www.zuckerjagdwurst.com/de/rezepte/veganer-sauerbraten-mit-spaetzle-und-rotkraut}{Zucker\&Jagdwurst: Veganer Sauerbraten}},
    calory = {405kcal | 64g Protein | 13g KH | 11g Fett},
    price = {2,35€},
    created = {10.06.2026},
    %rating = {1},
]
{Seitan-Sauerbraten}

%\graph{% pictures
%    big=pic/basics/07_seitan-sauerbraten
%}

\introduction{%
Ein kräftiger Seitanbraten als Basis für veganen Sauerbraten. Erst wird der Seitan gedämpft, dann über Nacht gekühlt und anschließend in einer Essigbeize mariniert. Also nichts für spontan, aber genau deswegen schmeckt es später auch nicht nach traurigem Gummiball.
}

\ingredients{
    \textbf{Seitanbraten} & \\
    500 g & Seitan-Fix\index{Tofu \& Fleischalternativen!Seitan}\\
    500 ml & Kalte Gemüsebrühe\\
    2 EL & Sojasauce\\
    2 EL & Senf\\
    2 EL & Hefeflocken\\
    2 TL & Paprikapulver\\
    2 TL & Knoblauchpulver\\
    1 TL & Zwiebelpulver\\
    1 TL & Thymian\\
    1 TL & Pfeffer\\
    1 TL & Salz\\
    1 Prise & Muskat\\
    1 Prise & Curry\\
    \textbf{Essigbeize} & \\
    500 ml & Wasser\\
    500 ml & Rotweinessig\\
    200 ml & Rotwein\\
    2 & Lorbeerblätter\\
    4 & Wacholderbeeren\\
    2 & Nelken\\
    3 & Zwiebeln\index{Gemüse!Zwiebel}\\
    \textbf{Zum Anbraten} & \\
    4 EL & Pflanzenöl
}

\preparation{%
    \step%
    Seitan-Fix mit Paprikapulver, Knoblauchpulver, Zwiebelpulver, Thymian, Pfeffer, Salz, Hefeflocken, Muskat und Curry in einer großen Schüssel vermischen.
    
    \step%
    Kalte Gemüsebrühe, Sojasauce und Senf zugeben und alles 5--10 Minuten kräftig durchkneten, bis eine kompakte, elastische Masse entsteht.
    
    \step%
    Die Masse zu einem großen Braten formen und 30 Minuten bei Zimmertemperatur ruhen lassen.
    
    \step%
    Einen großen Topf mit etwas Wasser füllen und einen Dämpfeinsatz einsetzen. Der Seitan sollte nicht im Wasser liegen.
    
    \step%
    Seitanbraten in den Dämpfeinsatz legen, Deckel aufsetzen und 75 Minuten dämpfen. Dabei regelmäßig kontrollieren, ob noch genug Wasser im Topf ist.
    
    \step%
    Den gedämpften Seitan vollständig abkühlen lassen und über Nacht in den Kühlschrank stellen.
    
    \step%
    Am nächsten Tag für die Beize Wasser, Rotweinessig und Rotwein vermischen. Lorbeerblätter, Wacholderbeeren, Nelken und halbierte Zwiebeln zugeben.
    
    \step%
    Den kalten Seitanbraten in die Beize legen und 24 Stunden marinieren lassen.
    
    \step%
    Danach den Seitan aus der Beize nehmen und gut abtropfen lassen. Die Beize aufheben, falls daraus später eine Sauerbraten-Sauce gekocht werden soll.
    
    \step%
    Pflanzenöl in einer großen Pfanne oder einem Bräter erhitzen und den Seitan von allen Seiten kräftig anbraten, bis er schön gebräunt ist.
}

\suggestion[Weiterverwenden]{%
    Für Sauerbraten den angebratenen Seitan in Scheiben schneiden oder mit zwei Gabeln grob zerreißen und anschließend in einer sauren Bratensauce ziehen lassen. Gerade das Zerreißen hilft, damit mehr Sauce in den Seitan kommt.
}

\hint{
    Seitan wird deutlich besser, wenn er nach dem Dämpfen wirklich komplett abkühlt und über Nacht ruht. Direkt warm weiterverarbeitet ist die Konsistenz oft deutlich weicher und weniger bratentauglich.
}

\end{recipe}